\chapter{Theoretical Background}
\label{ch:theory}

\section{Extended Hamilton Principle for Biofilm Mechanics}
\label{sec:theory_hamilton}

\subsection{State Variables and Kinematics}

Following Klempt et al.~\cite{Klempt2025ContinuumBiofilm} --- who extended the
single-species Hamilton-based growth model of~\cite{Klempt2024ContinuumGrowth} to
multiple interacting species --- the biofilm
state at a material point is described by two internal variable fields:
the \emph{volume fraction} $\phi_i \in [0,1]$ occupied by species $i$ ($i = 1,\dots,n$),
and the \emph{viability fraction} $\psi_i \in [0,1]$, representing the proportion of
living cells within species $i$.
The living volume fraction is then
\begin{equation}
  \bar{\phi}_i := \phi_i\,\psi_i, \qquad i = 1,\dots,n.
  \label{eq:living}
\end{equation}
An additional void fraction $\phi_0\in[0,1]$ acts as a numerical auxiliary variable
to enforce the holonomic incompressibility constraint~\cite{Klempt2025ContinuumBiofilm}:
\begin{equation}
  \sum_{l=0}^{n} \phi_l(t) = 1 \quad \forall\, t.
  \label{eq:constraint}
\end{equation}

\subsection{Variational Formulation}

The evolution equations are derived from the Extended Hamilton Principle~\cite{JunkerBalzani2021ExtendedHamilton}:
\begin{equation}
  \mathcal{H} = \int_\tau \bigl(\mathcal{G} + \mathcal{C} + \mathcal{D}\bigr)\,\mathrm{d}t
  \;\to\; \mathrm{stat}_{\boldsymbol{\xi},\gamma},
  \label{eq:hamilton_principle}
\end{equation}
where $\boldsymbol{\xi} = (\boldsymbol{\phi}, \boldsymbol{\psi})^\top$ collects
all state variables, $\gamma$ is a Lagrange multiplier enforcing the constraint
\eqref{eq:constraint}, $\mathcal{G}$ is the Helmholtz free energy, $\mathcal{C}$
the constraint functional, and $\mathcal{D}$ the dissipation.

The free energy density encodes nutrient-driven growth and antibiotic-induced killing:
\begin{equation}
  \Psi = -\tfrac{1}{2}\,c^*\,\bar{\bphi}^\top \bA\,\bar{\bphi}
        + \tfrac{1}{2}\,\alpha^*\,\bpsi^\top \mathbf{B}\,\bpsi,
  \label{eq:free_energy}
\end{equation}
where $c^*$ is the nutrient concentration, $\alpha^*$ the antibiotic concentration,
$\bA \in \Real^{n \times n}$ is the symmetric growth-interaction matrix, and
$\mathbf{B} = \mathrm{diag}(b_1,\dots,b_n)$ the antibiotic sensitivity matrix.
The volume constraint enters through
\begin{equation}
  \mathcal{C} = \gamma\!\left(\sum_{l=0}^{n} \phi_l - 1\right).
\end{equation}
The dissipation potential is~\cite{Klempt2025ContinuumBiofilm}
\begin{equation}
  \mathcal{D} = \tfrac{1}{2}\,\dot{\bar{\bphi}}^\top \boldsymbol{\eta}\,\dot{\bar{\bphi}}
              + \tfrac{1}{2}\,\dot{\bphi}^\top \boldsymbol{\eta}\,\dot{\bphi},
  \label{eq:dissipation}
\end{equation}
with diagonal viscosity matrix $\boldsymbol{\eta} = \mathrm{diag}(\eta_1,\dots,\eta_n)$,
$\eta_i = 1$ for all species (absorbed into the time scale).
The two-term structure is necessary: using only the first term would yield a
rank-deficient system because $\bar\phi_i=\phi_i\psi_i$ couples the two rate
variables; the second term $\tfrac12\dot\bphi^\top\boldsymbol\eta\,\dot\bphi$
restores full rank~\cite{Klempt2025ContinuumBiofilm}.

\subsection{Governing Equations}

Evaluating $\delta\mathcal{H} = 0$ with respect to $\delta\phi_i$, $\delta\psi_i$, and
$\delta\gamma$ yields:

\medskip\noindent
\textit{Volume fraction} ($i = 1,\dots,n$):
\begin{equation}
  0 = -c^*\,\psi_i\,\bigl(\bA\bar{\bphi}\bigr)_i
      + \eta_i\!\left(\dot\phi_i\,\psi_i^2 + \bar\phi_i\,\dot\psi_i + \dot\phi_i\right)
      + \gamma.
  \label{eq:gov_phi}
\end{equation}
\footnotetext{The dissipation terms follow from
$\partial\mathcal{D}/\partial\dot\phi_i
 = \eta_i\,\dot{\bar\phi}_i\,(\partial\bar\phi_i/\partial\phi_i)
 + \eta_i\,\dot\phi_i
 = \eta_i(\dot\phi_i\psi_i+\phi_i\dot\psi_i)\psi_i + \eta_i\dot\phi_i
 = \eta_i(\dot\phi_i\psi_i^2 + \bar\phi_i\dot\psi_i + \dot\phi_i)$,
using $\bar\phi_i=\phi_i\psi_i$ and the chain rule
$\partial\dot{\bar\phi}_i/\partial\dot\phi_i=\psi_i$.
The $+\dot\phi_i$ term without $\psi$-weighting arises from the second term
$\tfrac12\dot\bphi^\top\boldsymbol\eta\,\dot\bphi$ in~\eqref{eq:dissipation}.
The $\delta\psi_i$-equation follows analogously.}

\noindent\textit{Viability fraction} ($i = 1,\dots,n$):
\begin{equation}
  0 = -c^*\,\phi_i\,\bigl(\bA\bar{\bphi}\bigr)_i
      + \alpha^*\,b_i\,\psi_i
      + \eta_i\!\left(\dot\psi_i\,\phi_i^2 + \bar\phi_i\,\dot\phi_i\right)
      + \gamma.
  \label{eq:gov_psi}
\end{equation}

\noindent\textit{Void fraction and constraint:}
\begin{equation}
  0 = \gamma + \dot\phi_0, \qquad
  0 = \sum_{l=0}^{n}\phi_l - 1.
  \label{eq:gov_void}
\end{equation}

The system \eqref{eq:gov_phi}--\eqref{eq:gov_void} constitutes $2n+2$ coupled non-linear
equations, solved at each time step via Newton's method with implicit Euler discretisation,
the standard choice for the stiff system near the simplex boundary~\cite{Hairer1996ODE}.
In the numerical implementation a logarithmic barrier penalty $\Psi_{\mathrm{bar}}$ with
$K_p = \eta_i\times 10^{-4}$ enforces the box constraints $\phi_i,\psi_i\in(0,1)$, omitted
from \eqref{eq:gov_phi}--\eqref{eq:gov_void} for clarity.
Symmetry $A_{ij} = A_{ji}$ is a modelling assumption imposed
in~\cite{Klempt2025ContinuumBiofilm}: the free energy~\eqref{eq:free_energy} is
well-defined for any square $\bA$, but its antisymmetric part does not contribute
to the quadratic form and is therefore set to zero by convention, halving the
parameter count from $n^2$ to $n(n+1)/2$.
This choice preserves the gradient-flow structure of the dynamics and the
associated thermodynamic consistency guarantee~\cite{JunkerBalzani2021ExtendedHamilton},
in the tradition of non-equilibrium thermodynamics with internal variables and
extremum principles~\cite{Maugin1992NonequilibriumTD, Ziegler1963MaximumEntropy}.
Biologically, symmetry is a reasonable approximation when interactions are
mediated by \emph{diffusible metabolites shared through a common pool}: if
species $j$ secretes lactate that promotes species $i$ (cross-feeding,
$A_{ij}>0$), the same metabolite pool constrains $j$ in proportion, making
$A_{ij}\approx A_{ji}$ for syntrophic partners~\cite{Kolenbrander2010OralMultispecies}.
Genuinely asymmetric effects (directional toxin secretion, contact killing) are
suppressed and constitute the domain where the unconstrained gLV retains a small
predictive edge (Chapter~\ref{ch:dieckow}).

\paragraph{Scale gauge.}
The nutrient concentration $c^*$ enters the evolution equations only through the product
$c^*A_{ij}$, so its absolute value is non-identifiable: rescaling $c^*\mapsto\lambda c^*$ and
$\bA\mapsto\lambda^{-1}\bA$ leaves the dynamics invariant. We fix this gauge by setting
$c^*=25$ (cf.\ $c^*=100\,\mathrm{J\,m^{-3}}$ in the numerical experiments of
Klempt et al.~\cite{Klempt2025ContinuumBiofilm}), which renders the inferred
$A_{ij}=\mathcal{O}(1)$; $c^*$ therefore sets the
overall time scale rather than an independent degree of freedom. The Heine et al.\ assays
use no antibiotic, so $\alpha^*=0$ deactivates the decay term $\alpha^*b_i\psi_i$ and the
sensitivity matrix $\mathbf{B}$.

\paragraph{From composition to mechanics.}
The Extended Hamilton Principle above governs the \emph{composition} fields
$(\bphi,\bpsi)$; the biofilm's \emph{deformation} is the complementary mechanical degree
of freedom, governed by the same thermodynamic philosophy of a free energy rendered
stationary subject to constraints. Because the local biomass fraction sets the local
rate of biomass accretion, the composition acts as the driver of a stress-free
\emph{growth} of the continuum, which the standard multiplicative split
$\bF=\bF_{\!e}\bF_{\!g}$~\cite{Rodriguez1994Growth} separates from the stress-carrying
elastic deformation $\bF_{\!e}$; closing the description with a hyperelastic constitutive
law for $\bF_{\!e}$ then turns a depth-varying composition into a depth-varying
mechanical state. The compositional incompressibility constraint~\eqref{eq:constraint}
--- a mixture condition on the \emph{relative} fractions --- is distinct from, and
compatible with, the physical volume growth $J_g=\det\bF_{\!g}>1$ of accreting biomass.
The finite-element realisation of this composition\,$\rightarrow$\,growth\,$\rightarrow$\,stress
chain, and the residual stress it imprints at the implant interface, is developed in
Section~\ref{sec:integration_fem}.


\section{Replicator Equation and Compositional Dynamics}
\label{sec:theory_replicator}

The model integrated throughout this thesis is the full coupled $\bphi$--$\bpsi$
system~\eqref{eq:gov_phi}--\eqref{eq:gov_void}. It is instructive to relate it to the
classical \emph{replicator equation} of evolutionary game theory~\cite{Hofbauer1998EGT,TaylorJonker1978Replicator}:
\begin{equation}
  \dot{\phi}_i = \phi_i \Bigl[ \bigl(\bA\bphi\bigr)_i - \bphi^\top \bA \bphi \Bigr],
  \qquad \sum_i \phi_i = 1,
  \label{eq:replicator}
\end{equation}
in which the mean-fitness term $\bphi^\top\bA\bphi$ projects the growth onto the tangent
plane of the $n$-simplex, conserving composition exactly. For a symmetric $\bA$,
\eqref{eq:replicator} is a (generalised) gradient flow of the mean fitness
$W(\bphi)=\tfrac12\bphi^\top\bA\bphi$ with respect to the \emph{Shahshahani information
metric} $g_{ij}=\delta_{ij}/\phi_i$, and $W$ is a strict Lyapunov function---its monotone
increase (equivalently, the monotone decrease of the free energy $\Psi=-W$) is the
thermodynamic-consistency statement, the dissipation being non-negative by
construction~\cite{Hofbauer1998EGT, Shahshahani1979}.

\paragraph{Remark (mobility metric).}
The $\phi_i$-weighted form~\eqref{eq:replicator} corresponds specifically to the
Shahshahani (information-geometric) mobility. The continuum model's dissipation
potential~\eqref{eq:dissipation} uses a constant viscosity $\boldsymbol{\eta}$; its
well-mixed, $\psi_i\!\equiv\!1$ reduction is therefore the affine projection
$\dot\phi_i\propto(\bA\bphi)_i-\langle\bA\bphi\rangle$ rather than~\eqref{eq:replicator}
verbatim. The two coincide in sign structure and fixed points; \eqref{eq:replicator} is
retained as the canonical compositional reference because the inference operates on the
same symmetric $\bA$ and simplex constraint.

\paragraph{Relation to gLV.}
Equation~\eqref{eq:replicator} differs from the generalised Lotka-Volterra (gLV) model
\eqref{eq:glv} in two key respects: (i) the interaction matrix $\bA$ is constrained to
be \emph{symmetric} by the variational derivation, reducing the parameter count from
$n^2$ to $n(n+1)/2$; (ii) the right-hand side preserves $\sum_i \phi_i = 1$ exactly,
whereas gLV models require explicit normalisation and do not guarantee volume conservation.
These properties make the Hamilton replicator particularly suited for closed microbial
communities where total biovolume is constrained (e.g., in-vitro cultures of fixed total
cell concentration).

\paragraph{Fixed points and stability.}
The interior fixed points of~\eqref{eq:replicator} satisfy $(\bA\bphi^*)_i = \bphi^{*\top}\bA\bphi^*$
for all $i$ with $\phi^*_i > 0$.
For a positive definite $\bA$, the system has a unique interior fixed point, which is
a globally stable equilibrium on the simplex (Lyapunov function: $V = -\tfrac{1}{2}\bphi^\top\bA\bphi$).
In the dysbiotic condition, the interaction matrix inferred from Heine et al.'s data
exhibits a near-singular structure with large off-diagonal entries involving $\text{Pg}$,
consistent with the bistable dynamics observed experimentally.


\section{Bayesian Inference and TMCMC}
\label{sec:theory_tmcmc}

\subsection{Bayesian Parameter Identification}

Given observed species trajectories $\mathbf{y} = \{\mathbf{y}_t\}_{t=1}^{N_t}$
and a forward model $\mathcal{F}(\btheta)$ parameterised by the interaction vector
$\btheta \in \Real^{15}$, Bayesian inference seeks the posterior distribution:
\begin{equation}
  p(\btheta \mid \mathbf{y}) \;\propto\; \mathcal{L}(\mathbf{y} \mid \btheta)\,p(\btheta),
  \label{eq:bayes}
\end{equation}
where $\mathcal{L}(\mathbf{y} \mid \btheta)$ is the likelihood and $p(\btheta)$ the prior.
We adopt a Gaussian likelihood with isotropic noise:
\begin{equation}
  \mathcal{L}(\mathbf{y} \mid \btheta)
  = \prod_{t,i} \mathcal{N}\!\left(y_{t,i} \mid \hat\phi_i(t;\btheta),\,\sigma^2\right),
\end{equation}
where $\hat\phi_i(t;\btheta)$ is the ODE solution at time $t$.

\subsection{Transitional Markov Chain Monte Carlo}

TMCMC~\cite{ChingChen2007TMCMC, Betz2016TMCMC} constructs a sequence of intermediate
distributions that bridge the prior to the posterior via a tempering schedule
$0 = \beta_0 < \beta_1 < \cdots < \beta_J = 1$:
\begin{equation}
  p_j(\btheta) \propto \mathcal{L}(\mathbf{y} \mid \btheta)^{\beta_j}\,p(\btheta).
  \label{eq:tmcmc_intermediate}
\end{equation}
At each stage $j$, importance weights $w_j^{(k)} \propto \mathcal{L}^{\beta_j - \beta_{j-1}}$ are
computed, particles are resampled, and short Metropolis--Hastings chains are run to
diversify the ensemble.
The tempering exponent $\beta_j$ is chosen adaptively to maintain a target coefficient
of variation of the weights (typically $c_v = 1$), ensuring smooth progression from
the diffuse prior to the concentrated posterior.

\paragraph{GPU parallelisation.}
The dominant cost is the batched forward-model evaluation at each stage.
In this work, the ODE right-hand side is compiled to GPU-executable XLA code via
JAX~\cite{jax2018github}, and the entire particle ensemble (typically $N = 2{,}000$)
is evaluated in a single batched call using \texttt{jax.vmap}.
This yields a ${\sim}200\times$ speedup over the CPU baseline
(Section~\ref{sec:heine_gpu}), making TMCMC inference on the 15-dimensional
parameter vector practically feasible within minutes.

\paragraph{Iteratively narrowed priors.}
To address posterior multimodality in the dysbiotic condition, we employ an iterative
strategy: after an initial TMCMC run on a broad prior, the posterior mean and standard
deviation are used to construct a tighter prior for a subsequent run.
This two-pass approach concentrates sampling on the high-posterior-density region,
reducing the effective dimensionality of the problem and improving convergence.

\paragraph{Identifiability diagnostic.}
To distinguish data-informed parameters from prior-dominated ones, we use the
effective-dimensionality ratio
\begin{equation}
  r_i \;=\; \frac{\sigma_{\mathrm{post},i}}{\Delta_{\mathrm{prior},i}},
  \label{eq:effdim}
\end{equation}
the posterior standard deviation of $\theta_i$ relative to its prior width.
A value $r_i \ll 1$ indicates the datum tightly constrains $\theta_i$ (the posterior is
much narrower than the prior), whereas $r_i \to 1$ flags a prior-dominated, non-identified
parameter. Reporting $r_i$ per parameter makes the identifiability of the
$15$-dimensional inference explicit rather than assumed
(Chapter~\ref{ch:heine}: all $15$ entries satisfy $r_i < 0.43$, i.e.\ every parameter is
data-informed).


\section{Flux Balance Analysis and Metabolic Sign Priors}
\label{sec:theory_fba}

\subsection{Flux Balance Analysis}

Flux Balance Analysis (FBA) predicts the metabolic phenotype of a microorganism
under steady-state assumptions~\cite{OrthThiele2010WhatIsFBA, Heinken2023AGORA2}.
Given a stoichiometric matrix $S \in \Real^{m \times r}$ (metabolites $\times$ reactions),
FBA solves:
\begin{equation}
  \max_{\mathbf{v}}\; \mathbf{c}^\top\mathbf{v}
  \quad \text{s.t.} \quad S\mathbf{v} = \mathbf{0},\; \mathbf{lb} \le \mathbf{v} \le \mathbf{ub},
\end{equation}
where $\mathbf{c}$ is the objective vector (typically biomass production),
$\mathbf{v}$ the flux vector, and $\mathbf{lb}$, $\mathbf{ub}$ the reaction bounds.
Parsimonious FBA (pFBA)~\cite{Lewis2012pFBA} additionally minimises the total flux at the optimal growth
rate, yielding a biologically plausible flux distribution without unnecessary
enzyme activity.

\subsection{Community FBA via MICOM}

MICOM~\cite{Diener2020MICOM} extends FBA (computed here with COBRApy~\cite{Ebrahim2013COBRApy}) to multi-species consortia, generalising single-organism dynamic FBA~\cite{Mahadevan2002DynamicFBA}, by coupling
individual metabolic models through shared metabolite exchange fluxes and optimising a
cooperative tradeoff objective balancing community growth with individual growth.
For each species pair $(i, j)$, MICOM predicts the net exchange flux of shared metabolites,
from which the sign of ecological interaction --- mutualism ($+/+$), competition ($-/-$),
commensalism ($+/0$), or amensalism ($-/0$) --- can be inferred.

\subsection{Three-Layer Sign Prior}

For the Dieckow dataset (Chapter~\ref{ch:dieckow}), a three-layer sign prior
$\mathbf{S} \in \{-1, 0, +1\}^{n \times n}$ is constructed hierarchically:
\begin{enumerate}
  \item \textbf{Layer 1 (eHOMD experimental)}: Empirical co-occurrence and
    exclusion data from the expanded Human Oral Microbiome Database (eHOMD).
  \item \textbf{Layer 2 (AGORA2 pFBA)}: Parsimonious FBA predictions of
    pairwise metabolite exchange from 7,302 gut/oral reconstructions~\cite{Heinken2023AGORA2}.
  \item \textbf{Layer 3 (MICOM community FBA)}: Community-level cooperative
    tradeoff predictions resolving conflicts between Layers 1 and 2~\cite{Diener2020MICOM}.
\end{enumerate}
The sign prior is incorporated into the loss function as a half-space hinge-squared penalty:
\begin{equation}
  \mathcal{R}(\btheta; \mathbf{S}) = \lambda \sum_{i,j:\, S_{ij}\neq 0}
  \max\!\bigl(0,\; -S_{ij}\,A_{ij}\bigr)^2,
\end{equation}
penalising parameter values whose sign contradicts the metabolic prediction while
leaving sign-consistent entries unconstrained.

The prior uses the \emph{sign} of the metabolic signal exclusively, not its
magnitude.
This is not a modelling compromise but the principled choice given the current
state of constraint-based metabolic modelling.
Rath et al.\ \cite{Rath2024GEMReliability} evaluated semi-curated GEMs (including
AGORA-derived reconstructions) against in-vitro interaction data and found that
predicted interaction strengths are quantitatively unreliable (Spearman
$r < 0.3$), while the sign/direction of interactions retains practical predictive
signal.
A mechanistic derivation linking FBA exchange fluxes
[mmol\,gDW$^{-1}$\,h$^{-1}$] to dimensionless gLV interaction coefficients does
not exist in the published literature: the MacArthur consumer-resource reduction
to gLV (Marsland et al.\ \cite{Marsland2020MiCRM}) uses coarse-grained preference
parameters rather than genome-scale stoichiometry, and the gLV approximation
itself breaks down above a metabolic leakage threshold of $l_i \approx 0.2$
that lies within the biologically realistic range (Mustri et al.\
\cite{Mustri2025GLVABreakdown}).
Using sign information only is therefore the appropriate and maximally honest
use of the metabolic model output.
